% GENERATED FILE. Edit canonical lessons, then run npm run generate:books.
% canonical-source-hash: fnv1a64:ab157de5986b1574
% canonical-lessons: GU-C483-vedna, GU-C483-casko, GU-C483-khucvu, GU-C483-tarphadvu, GU-C483-dah

\chapter{Where It Hurts}
\label{ch:gu-where-it-hurts}
\hlchaptermodality{483}

\begin{chapteropening}
\textbf{By the end of this chapter:} \emph{I can describe different kinds of pain.}

Complete the last lesson of chapter 483: I can describe different kinds of pain.
\end{chapteropening}

\section[\texorpdfstring{\gu{વેદના} (vednā)}{વેદના (vednā)}]{\gu{વેદના} (vednā) — agony; deep pain}

\label{lesson:GU-C483-vedna}



\begin{quote}
Before the new one: say the Gujarati for flat on one's back, then the Gujarati for to bang.
\end{quote}

\subsection*{You'll want to know: \gu{વેદના}}
\textbf{\gu{વેદના}} — \emph{vednā} — \textquotedblleft{}agony; deep pain\textquotedblright{}.

\begin{quote}
\gu{દર્દીની} \gu{વેદના} \gu{જોઈને} \gu{નર્સ} \gu{ઉદાસ} \gu{થઈ}.
\end{quote}

\begin{quote}
\emph{dardīnī vednā joīne nars udās thaī.}
\end{quote}

\begin{quote}
Seeing the patient's agony, the nurse grew sad.
\end{quote}

A feminine noun.

\begin{cousinweb}
\textbf{\gu{વેદના}} (\emph{vednā}) is a tatsama, Sanskrit \emph{vedanā}, \textquotedblleft{}pain, torture, agony\textquotedblright{}, as Monier-Williams records it.
\end{cousinweb}

\begin{grammarlens}[title={What you've built}]
One more word for describing what you lived and felt.
\end{grammarlens}

\subsection*{Your turn}
\begin{itemize}
\raggedright
  \item \emph{Say it:} \emph{vednā}
  \item \emph{Say it:} \emph{vednā}, once more
  \item \emph{Say it:} \emph{dardīnī vednā joīne nars udās thaī.} — Seeing the patient's agony, the nurse grew sad.
  \item \emph{Recall:} say the Gujarati for flat on one's back, then the Gujarati for to bang, then say \emph{vednā} again
\end{itemize}

\begin{tcolorbox}[breakable,colback=teal!4,colframe=teal!35!black,title={Before you move on}]
What does \gu{વેદના} mean? (\textquotedblleft{}Agony; deep pain\textquotedblright{}.) Say it once more.
\end{tcolorbox}
\section[\texorpdfstring{\gu{ચસકો} (casko)}{ચસકો (casko)}]{\gu{ચસકો} (casko) — a twinge; a sudden sharp pain}

\label{lesson:GU-C483-casko}



\begin{quote}
Before the new one: say the Gujarati for to bang, then the Gujarati for agony.
\end{quote}

\subsection*{You'll want to know: \gu{ચસકો}}
\textbf{\gu{ચસકો}} — \emph{casko} — \textquotedblleft{}a twinge; a sudden sharp pain\textquotedblright{}.

\begin{quote}
\gu{સવારે} \gu{કમરમાં} \gu{ચસકો} \gu{આવ્યો}.
\end{quote}

\begin{quote}
\emph{savāre kamarmā̃ casko āvyo.}
\end{quote}

\begin{quote}
This morning I felt a twinge in my back.
\end{quote}

A masculine noun.

\begin{cousinweb}
\textbf{\gu{ચસકો}} (\emph{casko}) goes back to a reconstructed \emph{*cassakk-}, \textquotedblleft{}throb, twitch, sudden pain\textquotedblright{}; CDIAL lists the Gujarati \emph{casak}, \textquotedblleft{}acute pain\textquotedblright{} (CDIAL 4730).
\end{cousinweb}

\begin{grammarlens}[title={What you've built}]
One more word for describing what you lived and felt.
\end{grammarlens}

\subsection*{Your turn}
\begin{itemize}
\raggedright
  \item \emph{Say it:} \emph{casko}
  \item \emph{Say it:} \emph{casko}, once more
  \item \emph{Say it:} \emph{savāre kamarmā̃ casko āvyo.} — This morning I felt a twinge in my back.
  \item \emph{Recall:} say the Gujarati for to bang, then the Gujarati for agony, then say \emph{casko} again
\end{itemize}

\begin{tcolorbox}[breakable,colback=teal!4,colframe=teal!35!black,title={Before you move on}]
What does \gu{ચસકો} mean? (\textquotedblleft{}A twinge; a sudden sharp pain\textquotedblright{}.) Say it once more.
\end{tcolorbox}
\section[\texorpdfstring{\gu{ખૂંચવું} (khū̃cvũ)}{ખૂંચવું (khū̃cvũ)}]{\gu{ખૂંચવું} (khū̃cvũ) — to prick; to dig in}

\label{lesson:GU-C483-khucvu}



\begin{quote}
Before the new one: say the Gujarati for agony, then the Gujarati for a twinge.
\end{quote}

\subsection*{You'll want to know: \gu{ખૂંચવું}}
\textbf{\gu{ખૂંચવું}} — \emph{khū̃cvũ} — \textquotedblleft{}to prick; to dig in\textquotedblright{}.

\begin{quote}
\gu{નવાં} \gu{ચંપલ} \gu{પગમાં} \gu{ખૂંચે} \gu{છે}.
\end{quote}

\begin{quote}
\emph{navā̃ campal pagmā̃ khū̃ce che.}
\end{quote}

\begin{quote}
The new sandals dig into my feet.
\end{quote}

Past (it …): \gu{તે} \gu{સોય} \gu{આંગળીમાં} \gu{ખૂંચી}, \emph{te soy ā̃gḷīmā̃ khū̃cī}.

\begin{cousinweb}
\textbf{\gu{ખૂંચવું}} (\emph{khū̃cvũ}) goes back to a reconstructed \emph{*khucc-}, \textquotedblleft{}pierce, tear\textquotedblright{}; CDIAL lists the Gujarati verb, \textquotedblleft{}to pierce\textquotedblright{} (CDIAL 3890).
\end{cousinweb}

\begin{grammarlens}[title={What you've built}]
One more word for describing what you lived and felt.
\end{grammarlens}

\subsection*{Your turn}
\begin{itemize}
\raggedright
  \item \emph{Say it:} \emph{khū̃cvũ}
  \item \emph{Say it:} \emph{khū̃cvũ}, once more
  \item \emph{Say it:} \emph{navā̃ campal pagmā̃ khū̃ce che.} — The new sandals dig into my feet.
  \item \emph{Recall:} say the Gujarati for agony, then the Gujarati for a twinge, then say \emph{khū̃cvũ} again
\end{itemize}

\begin{tcolorbox}[breakable,colback=teal!4,colframe=teal!35!black,title={Before you move on}]
What does \gu{ખૂંચવું} mean? (\textquotedblleft{}To prick; to dig in\textquotedblright{}.) Say it once more.
\end{tcolorbox}
\section[\texorpdfstring{\gu{તરફડવું} (tarphaḍvũ)}{તરફડવું (tarphaḍvũ)}]{\gu{તરફડવું} (tarphaḍvũ) — to writhe; to toss in pain}

\label{lesson:GU-C483-tarphadvu}



\begin{quote}
Before the new one: say the Gujarati for a twinge, then the Gujarati for to prick.
\end{quote}

\subsection*{You'll want to know: \gu{તરફડવું}}
\textbf{\gu{તરફડવું}} — \emph{tarphaḍvũ} — \textquotedblleft{}to writhe; to toss in pain\textquotedblright{}.

\begin{quote}
\gu{પેટના} \gu{દુખાવાથી} \gu{દર્દી} \gu{આખી} \gu{રાત} \gu{તરફડતો} \gu{હતો}.
\end{quote}

\begin{quote}
\emph{peṭnā dukhāvāthī dardī ākhī rāt tarphaḍto hato.}
\end{quote}

\begin{quote}
The patient was writhing all night with a stomach ache.
\end{quote}

Past (he …): \gu{તે} \gu{દુખાવાથી} \gu{તરફડ્યો}, \emph{te dukhāvāthī tarphaḍyo}.

\begin{cousinweb}
\textbf{\gu{તરફડવું}} (\emph{tarphaḍvũ}) goes back to a reconstructed \emph{*taḍapphaḍ-}, \textquotedblleft{}agitate\textquotedblright{}; CDIAL lists the Gujarati verb, \textquotedblleft{}to toss about, be agitated\textquotedblright{} (CDIAL 5634).
\end{cousinweb}

\begin{grammarlens}[title={What you've built}]
One more word for describing what you lived and felt.
\end{grammarlens}

\subsection*{Your turn}
\begin{itemize}
\raggedright
  \item \emph{Say it:} \emph{tarphaḍvũ}
  \item \emph{Say it:} \emph{tarphaḍvũ}, once more
  \item \emph{Say it:} \emph{peṭnā dukhāvāthī dardī ākhī rāt tarphaḍto hato.} — The patient was writhing all night with a stomach ache.
  \item \emph{Recall:} say the Gujarati for a twinge, then the Gujarati for to prick, then say \emph{tarphaḍvũ} again
\end{itemize}

\begin{tcolorbox}[breakable,colback=teal!4,colframe=teal!35!black,title={Before you move on}]
What does \gu{તરફડવું} mean? (\textquotedblleft{}To writhe; to toss in pain\textquotedblright{}.) Say it once more.
\end{tcolorbox}
\section[\texorpdfstring{\gu{દાહ} (dāh)}{દાહ (dāh)}]{\gu{દાહ} (dāh) — a burning feeling}

\label{lesson:GU-C483-dah}



\begin{quote}
Before the new one: say the Gujarati for to prick, then the Gujarati for to writhe.
\end{quote}

\subsection*{You'll want to know: \gu{દાહ}}
\textbf{\gu{દાહ}} — \emph{dāh} — \textquotedblleft{}a burning feeling\textquotedblright{}.

\begin{quote}
\gu{તાવમાં} \gu{આખા} \gu{શરીરમાં} \gu{દાહ} \gu{થાય} \gu{છે}.
\end{quote}

\begin{quote}
\emph{tāvmā̃ ākhā śarīrmā̃ dāh thāy che.}
\end{quote}

\begin{quote}
With a fever the whole body feels as if it is burning.
\end{quote}

A masculine noun.

\begin{cousinweb}
\textbf{\gu{દાહ}} (\emph{dāh}) is a tatsama, Sanskrit \emph{dāha}, \textquotedblleft{}burning, heat\textquotedblright{}, as Monier-Williams records it.
\end{cousinweb}

\begin{grammarlens}[title={What you've built}]
One more word for describing what you lived and felt.
\end{grammarlens}

\subsection*{Your turn}
\begin{itemize}
\raggedright
  \item \emph{Say it:} \emph{dāh}
  \item \emph{Say it:} \emph{dāh}, once more
  \item \emph{Say it:} \emph{tāvmā̃ ākhā śarīrmā̃ dāh thāy che.} — With a fever the whole body feels as if it is burning.
  \item \emph{Recall:} say the Gujarati for to prick, then the Gujarati for to writhe, then say \emph{dāh} again
\end{itemize}

\begin{tcolorbox}[breakable,colback=teal!4,colframe=teal!35!black,title={Before you move on}]
What does \gu{દાહ} mean? (\textquotedblleft{}A burning feeling\textquotedblright{}.) Say it once more.
\end{tcolorbox}
